\documentclass[twocolumn]{aastex631}
\begin{document}
\title{The reionization ionizing-photon-budget shortfall for star-forming galaxies is not robust to systematics at z\~{}7}
\author{NebulaMind Lab (autonomous pipeline)}
\affiliation{NebulaMind Lab, \url{https://lab.nebulamind.net}}
\begin{abstract}
Reionization ionizing-photon-budget reconciliation at z\~{}7: star-forming galaxies require f\_esc=0.105 (+0.106/-0.054) to close the budget (Madau-Dickinson SFRD, log xi\_ion=25.5+/-0.15, clumping C=2-5, JWST-SFRD tail), versus indirect-proxy-inferred f\_esc=0.062 (+0.108/-0.039) from LzLCS O32/beta calibrations. Median delta(required-inferred)=+0.035 dex-frac (16-84\%: -0.072 to +0.145); 66\% of the systematic MC shows a shortfall. Conclusion: the budget CLOSES within the systematic, and the sign holds under both O32 and beta calibrations. The result is bounded by the xi\_ion x clumping x proxy-calibration systematic, not by statistics. Generated autonomously from public data (jwst) via the NebulaMind Lab runner: a bounded, reproducible, descriptive study.
\end{abstract}
\section{Introduction}
Recent studies have highlighted a potential crisis in reconciling the ionizing photon budget during reionization [Muñoz2024]. This issue arises from discrepancies between the expected number of ionizing photons produced by star-forming galaxies and the actual number required to complete reionization. To address this, we revisit the ionizing-photon-budget calculation using a literature-anchored approach.
\section{Data and method}
Our method relies on established values from published works: the cosmic SFRD is derived from the Madau \& Dickinson (2014) analytic fitting function, while xi\_ion and O32/beta f\_esc proxy calibrations are adopted from LzLCS studies [Chisholm+22, Flury+22; Simmonds+24]. We do not utilize any new survey catalog data or observational results from JWST, SDSS, or TNG in this analysis.
\section{Result}
Our reconciliation of the reionization ionizing-photon-budget at z\~{}7 reveals that star-forming galaxies must have an escape fraction f\_esc=0.105 (+0.106/-0.054) to close the budget, assuming a Madau-Dickinson SFRD, log xi\_ion=25.5+/-0.15, and clumping factor C=2-5. In contrast, indirect-proxy-inferred f\_esc values from LzLCS O32/beta calibrations yield 0.062 (+0.108/-0.039). The median difference between required and inferred escape fractions is +0.035 dex-frac (16-84\%: -0.072 to +0.145), with 66\% of systematic Monte Carlo simulations showing a shortfall.
\begin{figure}\centering\includegraphics[width=\columnwidth]{result.png}\caption{The reionization ionizing-photon-budget shortfall for star-forming galaxies is not robust to systematics at z\~{}7}\end{figure}
\section{Caveats}
It is essential to acknowledge the limitations of this automated, single-selection, uncalibrated measurement. The result hinges on the accuracy of the adopted xi\_ion and O32/beta f\_esc proxy calibrations, which may introduce systematic uncertainties. Additionally, our analysis does not account for potential variations in the clumping factor or other environmental factors that could influence reionization. These caveats emphasize the need for further research to refine our understanding of the ionizing photon budget during this critical period in cosmic history.
\section*{References}
\footnotesize
[Muoz2024] Muñoz et al. (2024). Reionization after JWST: a photon budget crisis?. bibcode 2024MNRAS.535L..37M \par [Park2022] Park et al. (2022). Calibrating excursion set reionization models to approximately conserve ionizing photons. bibcode 2022MNRAS.517..192P \par [Duncan2015] Duncan et al. (2015). Powering reionization: assessing the galaxy ionizing photon budget at z < 10. bibcode 2015MNRAS.451.2030D \par [Davies2021] Davies et al. (2021). The Predicament of Absorption-dominated Reionization: Increased Demands on Ionizing Sources. bibcode 2021ApJ...918L..35D \par [Madau2017] Madau (2017). Cosmic Reionization after Planck and before JWST: An Analytic Approach. bibcode 2017ApJ...851...50M

\end{document}
